\documentclass[10pt,a4paper]{article}

% --- Packages -------------------------------------------------------------
\usepackage[utf8]{inputenc}
\usepackage[T1]{fontenc}
\usepackage[margin=1.8cm]{geometry}
\usepackage{graphicx}
\usepackage{booktabs}
\usepackage{amsmath}
\usepackage{hyperref}
\usepackage{caption}
\usepackage{titlesec}
\captionsetup{font=small,labelfont=bf,skip=2pt}

% Tight spacing
\setlength{\parskip}{1.5pt}
\setlength{\parindent}{0pt}
\titlespacing*{\section}{0pt}{4pt}{2pt}
\titlespacing*{\subsection}{0pt}{3pt}{1pt}
\titlespacing*{\subsubsection}{0pt}{2pt}{1pt}
\titleformat*{\section}{\normalsize\bfseries}
\titleformat*{\subsection}{\small\bfseries}

% Compact title
\makeatletter
\renewcommand{\maketitle}{%
    \begingroup
    \begin{center}
        {\large\bfseries \@title \par}
        \vspace{2pt}
        {\small Marco Cattaneo Vittone \quad Mario Gonzalez del Camino \quad Mohammed Ali \par}
        \vspace{1pt}
        {\footnotesize University of Bern --- April 2026 \par}
    \end{center}
    \vspace{-4pt}
    \endgroup
}
\makeatother

\title{Deep Learning for Accelerated MRI Reconstruction with Explainability Analysis}

\begin{document}
\maketitle

% =========================================================================
\section*{Abstract}
\small
\noindent
This project implements and evaluates a deep learning model for accelerated
MRI reconstruction, establishes classical baselines, validates on a structurally
distinct second dataset, and extends the analysis with multiple explainability
techniques. A U-Net is trained to reconstruct missing axial slices from their
two acquired neighbours (2.5D interpolation) under $\times2$ acceleration on
the IXI dataset, reaching \textbf{36.07\,dB PSNR} and \textbf{0.946 SSIM} ---
\textbf{+2.86\,dB} over linear interpolation. A controlled ablation
(PlainCNN, identical architecture without skip connections) collapses to
33.06\,dB --- \emph{below} the linear baseline --- identifying skip
connections as the essential ingredient. Cross-dataset evaluation on
BraTS2020 yields a foreground MAE of $0.00754$, confirming generalisation to
tumour-bearing, skull-stripped data. Four explainability methods consistently
show the model uses both input slices symmetrically (Left/Right
ratio $\approx 1.01$--$1.06$).
\normalsize

% =========================================================================
\section{Introduction}
Magnetic resonance imaging is slow because each axial slice is acquired
sequentially. Long scans reduce patient comfort, raise costs, and limit
scanner throughput. A natural acceleration strategy is to acquire only every
$k$-th slice and reconstruct the missing ones in software. We study $k{=}2$
(50\,\% missing) and frame reconstruction as image-to-image regression:
predict the missing slice $\hat{y}$ from its acquired neighbours
$(x_{\text{left}}, x_{\text{right}})$.
This work (i) implements the full preprocessing--training--evaluation
pipeline; (ii) establishes classical interpolation baselines; (iii) trains a
U-Net and a controlled \textbf{PlainCNN ablation}; (iv) validates on
BraTS2020; and (v) applies four complementary explainability methods.

% =========================================================================
\section{Datasets}
\textbf{IXI} (primary). About 600 brain MRI scans from healthy volunteers
acquired at three London hospitals on different scanners and field strengths.
Only T1-weighted MPRAGE sequences are used; 581 scans are retained after
quality filtering, producing \textbf{42\,816 input--target pairs} split
70/15/15 into train/val/test (seed 42).
\textbf{BraTS2020} (cross-dataset). 57\,195 pre-extracted axial HDF5 slices,
T1 channel, skull-stripped, including tumour-bearing scans. Same 2.5D
construction $\Rightarrow$ \textbf{28\,413 pairs}.

% =========================================================================
\section{Methodology}

\subsection{Preprocessing and pair construction}
Each volume is loaded from NIfTI, resized to $256\times256$ (bicubic with
anti-aliasing) and intensity-normalised to $[0,1]$ via per-volume min--max
scaling. Under $\times2$ acceleration every second slice is discarded;
for each missing slice the model receives the two acquired neighbours as a
two-channel input and predicts the intermediate slice.

\subsection{Baselines}
\textbf{Linear interpolation} (pixel-wise average of neighbours, weighted by
fractional position) and \textbf{cubic spline} (fit through all acquired
slices along depth, evaluated at missing positions). Both evaluated on the
same test set.

\subsection{Model and ablation}
The reconstruction model is a \textbf{U-Net} with four encoder stages
(channels $[32,64,128,256]$), 512-channel bottleneck, and a mirrored decoder
with bilinear upsampling and concatenated skip connections at each
resolution (\textbf{7.76\,M parameters}; Sigmoid output $\in[0,1]$).
\textbf{PlainCNN} is identical (same depth, widths, blocks, parameter count)
but with the skip connections \emph{removed}, isolating their architectural
contribution.

\subsection{Training}
Hybrid loss
$\mathcal{L} = 0.8\,\lVert \hat{y}-y\rVert_{1} + 0.2\,(1-\mathrm{SSIM}(\hat{y},y))$.
Adam, lr $10^{-3}$, weight decay $10^{-5}$, \texttt{ReduceLROnPlateau}
(patience 5, factor 0.5), batch size 8, up to 20 epochs (best-validation
checkpoint kept), RTX~5070~Ti.

\subsection{Metrics and explainability}
Reported: PSNR (dB), SSIM, MAE on the held-out test split. For BraTS we add
\textbf{FG-MAE} (MAE on foreground voxels, ground-truth $>0.02$) because
skull-stripped volumes are $\sim$70--80\,\% zeros that trivialise raw PSNR.
Four explainability methods are run on 50 uniformly-sampled test cases:
(i) \textbf{Grad-CAM} on the final decoder block;
(ii) \textbf{Integrated Gradients} (50 path steps, zero baseline);
(iii) \textbf{slice-channel ablation} (zero one input slice, measure output
change);
(iv) \textbf{occlusion sensitivity} ($16{\times}16$ patch slid over each
input). Predictive uncertainty is estimated by attaching a forward-hook
dropout ($p{=}0.3$) to the bottleneck and running 20 stochastic passes
(MC-Dropout).

% =========================================================================
\section{Results}

\subsection{Reconstruction performance and ablation}
The U-Net beats both classical baselines by \textbf{+2.86\,dB} PSNR and
$+0.04$ SSIM (Table~\ref{tab:ixi}). Crucially, the PlainCNN ablation
(33.06\,dB) is \emph{worse than linear interpolation} (33.21\,dB): a deep
network without skip connections cannot match a four-line classical
baseline. Skip connections account for a $\sim$3\,dB margin on their own.
The improvement is most pronounced in mid-volume slices where tissue is
densest and interpolation is hardest; easy boundary slices are reconstructed
near-perfectly by all methods.

\begin{table}[h]
    \centering
    \footnotesize
    \begin{tabular}{lccc}
        \toprule
        Method                & PSNR (dB)        & SSIM              & MAE      \\
        \midrule
        Cubic spline          & 33.19            & 0.9048            & ---      \\
        Linear interpolation  & 33.21            & 0.9066            & ---      \\
        PlainCNN (no skips)   & 33.06            & ---               & ---      \\
        \textbf{U-Net (ours)} & \textbf{36.07}   & \textbf{0.9461}   & \textbf{0.0093} \\
        \bottomrule
    \end{tabular}
    \caption{IXI test-set reconstruction ($n{=}6\,422$ pairs). U-Net std:
    $\pm6.5$\,dB PSNR, $\pm0.029$ SSIM.}
    \label{tab:ixi}
\end{table}

\subsection{Cross-dataset generalisation}
The same architecture, retrained on BraTS pairs, attains
\textbf{FG-MAE 0.00754} on tissue voxels (raw PSNR 58.65\,dB but inflated by
the dominant skull-stripped background). The model generalises to a
structurally different dataset (tumour-bearing, skull-stripped) without
architectural changes.

\subsection{Explainability}
Three independent input-attribution methods (IG, channel ablation,
occlusion) yield Left/Right ratios within $\sim$6\,\% of unity
(Table~\ref{tab:expl}), strongly confirming a genuinely \emph{bidirectional}
interpolation. Channel ablation produces large output changes ($\sim 0.08$)
while patch occlusion produces small changes ($\sim 3\!\times\!10^{-4}$),
indicating reliance on \emph{global} slice structure rather than local
features --- consistent with anatomical interpolation. Grad-CAM scales
coherently with reconstruction difficulty: hard cases produce diffuse
attention, medium cases focus on the skull boundary and tissue edges, easy
cases show near-zero activation.

\begin{table}[h]
    \centering
    \footnotesize
    \begin{tabular}{lcc}
        \toprule
        Metric                                & Value                          & Interpretation \\
        \midrule
        IG L/R ratio                          & \textbf{1.0093}                & symmetric inputs \\
        Slice contribution L/R ratio          & \textbf{1.059}                 & symmetric ($\sim$6\,\%) \\
        Occlusion sensitivity L/R ratio       & \textbf{1.028}                 & no local region dominates \\
        MC variance (20 passes)               & $7.44\!\times\!10^{-16}$       & confident on test set \\
        \bottomrule
    \end{tabular}
    \caption{Explainability summary across 50 test samples. Per-method
    means $\pm$ stds in Appendix~\ref{app:expl}.}
    \label{tab:expl}
\end{table}

% =========================================================================
\section{Discussion, limitations and future work}

The PlainCNN ablation is the strongest finding: a deep CNN without skip
connections is \emph{worse} than 4-tap linear interpolation. Without skips,
all spatial information must survive a $16{\times}16$ bottleneck --- too
coarse for $256{\times}256$ reconstruction. The $+3$\,dB the skip
connections add is therefore the difference between \emph{useful} and
\emph{useless}, not an incremental gain. Cross-dataset transfer to BraTS
works without architectural changes, but PSNR alone is misleading on
skull-stripped data; FG-MAE is the honest measure. The explainability
findings are reassuring: a well-behaved interpolation model \emph{should}
treat its two neighbours symmetrically, and three independent methods
confirm ours does.

\textbf{Limitations.} (i) IXI is healthy-only; the model has not seen
tumours, lesions or hemorrhages on IXI (BraTS partially addresses this
but does not replace a clinical study). (ii) PSNR/SSIM measure average
quality and underweight small focal abnormalities. (iii) No radiologist
reading study has been performed.

\textbf{Future work.} \emph{Adaptive acquisition}: acquire every second
slice, run anomaly detection, reconstruct only if no findings; otherwise
complete the scan. \emph{Uncertainty-gated safety}: MC variance is
near-zero on in-distribution data, so any threshold becomes meaningful only
on out-of-distribution inputs. \emph{Pathological validation} with paired
healthy/pathological scans and lesion annotations. \emph{DTI integration}
correlating reconstruction errors with white-matter tract locations.

% =========================================================================
% APPENDIX (does not count toward 3-page limit)
% =========================================================================
\newpage
\appendix
\section*{\centering Appendix}
\addcontentsline{toc}{section}{Appendix}

\section{Reproducibility}
\label{app:repro}
All code, configuration, trained weights, and per-sample explainability
outputs are in the project repository. Three commands reproduce the IXI
results; a fourth and fifth reproduce the ablation and BraTS results:
\begin{verbatim}
python -m src.preprocess
python -m src.train --model unet     --data ixi
python -m src.train --model plaincnn --data ixi   # ablation
python -m src.preprocess_brats
python -m src.train --model unet     --data brats
python -m src.explainability
\end{verbatim}

\section{Implementation details}
\label{app:impl}
\textbf{Memory.} Preprocessed arrays are stored on disk as three-phase
NumPy memmaps (count, allocate, fill) to avoid loading the full 42\,816-pair
($\approx$22\,GB) dataset into RAM.
\textbf{U-Net.} Each stage uses two
\texttt{Conv3$\times$3 $\to$ BatchNorm $\to$ ReLU} blocks; downsampling is
\texttt{MaxPool2d(2)}; upsampling is bilinear (\texttt{F.interpolate}).
\texttt{Sigmoid} on the last layer constrains output to $[0,1]$.
\textbf{Loss weights.} $\alpha{=}0.8$ ($L_{1}$) was selected on a
preliminary sweep; $\alpha\in\{0.5,0.7,0.8,0.9\}$ produced
PSNR within $\pm0.4$\,dB of the chosen value, with $0.8$ giving the best
SSIM/PSNR trade-off.
\textbf{FG-MAE definition.} For BraTS we mask
ground-truth voxels with $y > 0.02$ and report
$\mathrm{MAE}_{\text{fg}} = \tfrac{1}{|M|}\sum_{(i,j)\in M}|\hat{y}_{ij}-y_{ij}|$.

\section{Per-method explainability statistics}
\label{app:expl}
\begin{table}[h]
    \centering
    \footnotesize
    \begin{tabular}{lcc}
        \toprule
        Metric                                & Mean                            & Std \\
        \midrule
        IG attribution -- Left slice          & $9.20\!\times\!10^{-7}$         & $5.30\!\times\!10^{-7}$ \\
        IG attribution -- Right slice         & $9.11\!\times\!10^{-7}$         & $5.26\!\times\!10^{-7}$ \\
        Mean Grad-CAM activation              & $0.00858$                       & $0.00581$ \\
        Slice contribution -- Left            & $0.0808$                        & $0.0485$ \\
        Slice contribution -- Right           & $0.0763$                        & $0.0456$ \\
        Occlusion sensitivity -- Left         & $0.00033$                       & $0.00021$ \\
        Occlusion sensitivity -- Right        & $0.00031$                       & $0.00020$ \\
        MC dropout pixel variance             & $7.44\!\times\!10^{-16}$        & $5.24\!\times\!10^{-16}$ \\
        \bottomrule
    \end{tabular}
    \caption{Per-method explainability statistics across 50 IXI test
    samples (extended view of Table~\ref{tab:expl}).}
\end{table}

\section{Notes on metric inflation in skull-stripped data}
\label{app:fgmae}
A skull-stripped BraTS slice contains 70--80\,\% exactly-zero background.
Both the prediction and the ground truth equal zero in those regions, so
they contribute zero squared error. PSNR is dominated by the few non-zero
foreground pixels but normalised by the total pixel count, producing values
($\approx 58$\,dB) that look near-perfect. FG-MAE strips this artefact by
restricting the average to foreground voxels.

\section{References}
\label{app:refs}
\footnotesize
\begin{enumerate}
    \item Ronneberger, O., Fischer, P., \& Brox, T. (2015). U-Net:
          Convolutional Networks for Biomedical Image Segmentation.
          \emph{MICCAI}.
    \item Wang, Z., Bovik, A. C., Sheikh, H. R., \& Simoncelli, E. P.
          (2004). Image Quality Assessment: From Error Visibility to
          Structural Similarity. \emph{IEEE TIP}.
    \item Selvaraju, R. R., et al.\ (2017). Grad-CAM: Visual Explanations
          from Deep Networks via Gradient-based Localization. \emph{ICCV}.
    \item Sundararajan, M., Taly, A., \& Yan, Q. (2017). Axiomatic
          Attribution for Deep Networks. \emph{ICML}.
    \item Zeiler, M. D., \& Fergus, R. (2014). Visualizing and
          Understanding Convolutional Networks. \emph{ECCV}.
    \item Gal, Y., \& Ghahramani, Z. (2016). Dropout as a Bayesian
          Approximation. \emph{ICML}.
    \item IXI Dataset, Imperial College London,
          \url{https://brain-development.org/ixi-dataset/}.
    \item Bakas, S., et al.\ (2017--2020). BraTS Challenge.
          \emph{Nature Scientific Data} and subsequent.
    \item Paszke, A., et al.\ (2019). PyTorch: An Imperative Style,
          High-Performance Deep Learning Library. \emph{NeurIPS}.
\end{enumerate}

\end{document}
